\documentclass[11pt]{article}
\usepackage[margin=0.8in]{geometry}
\usepackage{amsmath,amssymb}
\usepackage{xcolor,listings}
\usepackage{hyperref}
\hypersetup{colorlinks=true,linkcolor=blue}
\lstset{language=Python,basicstyle=\ttfamily\small,backgroundcolor=\color{black!3},frame=single,rulecolor=\color{black!15},columns=fullflexible,keepspaces=true,showstringspaces=false,aboveskip=8pt,belowskip=8pt}
\setlength{\parindent}{0pt}
\setlength{\parskip}{6pt}
\newcommand{\ket}[1]{\lvert #1\rangle}
\newcommand{\vfour}[4]{\begin{pmatrix}#1\\#2\\#3\\#4\end{pmatrix}}
\newcommand{\mtwo}[4]{\begin{pmatrix}#1&#2\\#3&#4\end{pmatrix}}
\begin{document}
\begin{center}
{\LARGE 09: Gates acting on an MPS}\\[5pt]
{\large Full state $\longrightarrow$ MPS mathematics $\longrightarrow$ MPS code}
\end{center}
We apply Hadamard to the first spin, then CNOT to both spins. These are unitary gate operations; we have not yet chosen a Hamiltonian for time evolution.

\section*{1. Full-state evolution}
Use the basis $\ket{00},\ket{01},\ket{10},\ket{11}$, with $\ket{0}=\ket{\uparrow}$ and $\ket{1}=\ket{\downarrow}$.
\[
\ket{\psi_0}=\ket{00}=\vfour{1}{0}{0}{0},\qquad
H=\frac1{\sqrt2}\mtwo{1}{1}{1}{-1},\qquad
G=\begin{pmatrix}1&0&0&0\\0&1&0&0\\0&0&0&1\\0&0&1&0\end{pmatrix}.
\]
First apply $H\otimes I$:
\[
\ket{\psi_H}
=\frac1{\sqrt2}\begin{pmatrix}1&0&1&0\\0&1&0&1\\1&0&-1&0\\0&1&0&-1\end{pmatrix}
\vfour{1}{0}{0}{0}
=\frac1{\sqrt2}\vfour{1}{0}{1}{0}
=\frac{\ket{00}+\ket{10}}{\sqrt2}.
\]
Then apply CNOT:
\[
\ket{\psi_G}
=\frac1{\sqrt2}\begin{pmatrix}1&0&0&0\\0&1&0&0\\0&0&0&1\\0&0&1&0\end{pmatrix}
\vfour{1}{0}{1}{0}
=\frac1{\sqrt2}\vfour{1}{0}{0}{1}
=\frac{\ket{00}+\ket{11}}{\sqrt2}.
\]
CNOT flips the second spin when the first is $1$. The result is a Bell state.

\begin{lstlisting}
import numpy as np

up = np.array([1, 0], dtype=complex)
state = np.kron(up, up)
H = np.array([[1, 1], [1, -1]], dtype=complex) / np.sqrt(2)
CNOT = np.array([[1, 0, 0, 0],
                 [0, 1, 0, 0],
                 [0, 0, 0, 1],
                 [0, 0, 1, 0]], dtype=complex)
identity = np.eye(2, dtype=complex)

state_after_H = np.kron(H, identity) @ state
state_after_CNOT = CNOT @ state_after_H
\end{lstlisting}

\newpage
\section*{2. MPS evolution with ordinary matrices}
The coefficient matrix stores the state amplitudes:
\[
\ket{\psi}=C_{00}\ket{00}+C_{01}\ket{01}+C_{10}\ket{10}+C_{11}\ket{11},
\qquad C=\mtwo{C_{00}}{C_{01}}{C_{10}}{C_{11}}.
\]
It is not an operator acting on $\ket{00}$. For two spins, an MPS is a factorization
\[
C=A^{[1]}A^{[2]},\qquad
C_{ij}=\sum_\alpha A^{[1]}_{i\alpha}A^{[2]}_{\alpha j}.
\]
\textbf{Initial MPS.} The initial coefficient matrix factors directly:
\[
C_0=\mtwo{1}{0}{0}{0}
=\underbrace{\begin{pmatrix}1\\0\end{pmatrix}}_{A_0^{[1]}}
\underbrace{\begin{pmatrix}1&0\end{pmatrix}}_{A_0^{[2]}}.
\]
\textbf{Hadamard.} Update the first factor:
\[
A_H^{[1]}=HA_0^{[1]}=\frac1{\sqrt2}\begin{pmatrix}1\\1\end{pmatrix},
\qquad
C_H=A_H^{[1]}A_0^{[2]}=\frac1{\sqrt2}\mtwo{1}{0}{1}{0}.
\]
\textbf{CNOT.} Join the factors, flatten row by row, apply $G$, and reshape:
\[
\underbrace{\frac1{\sqrt2}\mtwo{1}{0}{1}{0}}_{C_H}
\xrightarrow{\text{flatten}}
\frac1{\sqrt2}\vfour{1}{0}{1}{0}
\xrightarrow{G}
\frac1{\sqrt2}\vfour{1}{0}{0}{1}
\xrightarrow{\text{reshape}}
\underbrace{\frac1{\sqrt2}\mtwo{1}{0}{0}{1}}_{C_G}.
\]
\textbf{Restore the MPS.} Factor the updated matrix:
\[
C_G=USV^\dagger,\qquad
A_G^{[1]}=U,\qquad A_G^{[2]}=SV^\dagger.
\]
One valid SVD gives
\[
U=V^\dagger=I,\qquad S=\frac1{\sqrt2}I,
\qquad
\boxed{A_G^{[1]}=\mtwo{1}{0}{0}{1},\quad
A_G^{[2]}=\frac1{\sqrt2}\mtwo{1}{0}{0}{1}.}
\]
The shared dimension has grown from $1$ to $2$. These new factors are the final MPS. An SVD routine may return different factors with the same product.

\textbf{Connection to the code.} Above, we flattened the coefficients to use the $4\times4$ gate. The code instead reshapes the gate and keeps the coefficient tensor's axes. Both routes give the same $C_G$ before SVD.

For two spins, $C_G$ holds the entire state's amplitudes. In a longer chain, the same procedure joins, updates, and splits only the two neighboring tensors touched by the gate; the full state need not be reconstructed.

\newpage
\section*{3. MPS code: initial tensors and Hadamard}
The code stores each tensor as
\[
(\text{left bond},\ \text{physical spin},\ \text{right bond}).
\]
For the initial product state:
\begin{lstlisting}
A1 = up.reshape(1, 2, 1).copy()
A2 = up.reshape(1, 2, 1).copy()
\end{lstlisting}
The entries are unchanged:
\[
\texttt{up[0]}=1,\quad\texttt{up[1]}=0
\quad\longrightarrow\quad
\texttt{A1[0,0,0]}=1,\quad\texttt{A1[0,1,0]}=0.
\]
Only the middle index can vary. The outer axes are placeholders: there are no spins beyond the two ends. Adding an axis of size one adds no coefficients.

Omitting those outer axes recovers our matrices:
\[
\texttt{A1[0,:,:]}=\begin{pmatrix}1\\0\end{pmatrix},
\qquad
\texttt{A2[:,:,0]}=\begin{pmatrix}1&0\end{pmatrix}.
\]
The shared dimension is one because the initial state is a product. After entanglement is created, it becomes two:
\[
\begin{array}{c|cc}
 & A^{[1]} & A^{[2]}\\\hline
\text{Initial shapes} &(1,2,1)&(1,2,1)\\
\text{Final shapes} &(1,2,2)&(2,2,1)
\end{array}
\]
An entangled state's tensors come from factoring the joint coefficients, not from reshaping two independent spin vectors.

\subsection*{Apply Hadamard}
\[
A_H^{[1]}=HA_0^{[1]}
\quad\longleftrightarrow\quad
(A_H^{[1]})_{aib}=\sum_j H_{ij}(A_0^{[1]})_{ajb}.
\]
\begin{lstlisting}
A1 = np.einsum("ij,ajb->aib", H, A1)
\end{lstlisting}
Here $a,b$ are bond indices; $i,j$ are spin indices. The gate acts on the physical index and leaves the bond axes in place.

An optional check reconstructs the state:
\begin{lstlisting}
mps_after_H = np.tensordot(
    A1, A2, axes=([-1], [0])
).reshape(-1)
print(np.linalg.norm(mps_after_H - state_after_H))
\end{lstlisting}
\[
A_H^{[1]}A_0^{[2]}=C_H
\xrightarrow{\text{flatten}}
\frac1{\sqrt2}\vfour{1}{0}{1}{0}.
\]

\newpage
\section*{4. MPS code: join and apply CNOT}
\subsection*{Join the two tensors}
\begin{lstlisting}
theta = np.tensordot(A1, A2, axes=([-1], [0]))
\end{lstlisting}
This sums their shared bond:
\[
(\Theta_H)_{aijc}=\sum_b (A_H^{[1]})_{aib}(A_0^{[2]})_{bjc}.
\]
For our two-spin state,
\[
\boxed{(\Theta_H)_{0ij0}=(C_H)_{ij}},\qquad
C_H=\frac1{\sqrt2}\mtwo{1}{0}{1}{0}.
\]
Thus \texttt{theta} is $C_H$ with its outer axes retained: shape $(1,2,2,1)$. This repeats the optional check's contraction, but does not flatten it.

\subsection*{Reshape the gate, then apply it}
The gate's row labels are $00,01,10,11$, and so are its column labels. Separate each label into its two spins:
\[
\underbrace{G_{2i+j,\,2p+q}}_{\text{matrix entry}}
\quad\xrightarrow{\text{reshape}}\quad
\underbrace{\texttt{gate[i,j,p,q]}}_{\text{same entry}}.
\]
\begin{lstlisting}
gate = CNOT.reshape(2, 2, 2, 2)
theta = np.einsum("ijpq,apqc->aijc", gate, theta)
\end{lstlisting}
The input spins are $p,q$; the output spins are $i,j$. The multiplication is
\[
(\Theta_G)_{aijc}
=\sum_{p,q}G_{2i+j,\,2p+q}(\Theta_H)_{apqc}.
\]
The code reuses the name \texttt{theta}: it now holds $\Theta_G$. The result is
\[
\underbrace{\bigl[(\Theta_H)_{0ij0}\bigr]_{ij}}_{C_H}
=\frac1{\sqrt2}\mtwo{1}{0}{1}{0}
\quad\xrightarrow{\text{CNOT}}\quad
\underbrace{\bigl[(\Theta_G)_{0ij0}\bigr]_{ij}}_{C_G}
=\frac1{\sqrt2}\mtwo{1}{0}{0}{1}.
\]
\subsection*{The equivalent vector route}
For this two-spin example, the two gate lines above could instead be:
\begin{lstlisting}
vector = theta.reshape(4)
vector = CNOT @ vector
theta = vector.reshape(1, 2, 2, 1)
\end{lstlisting}
These are alternatives, not consecutive steps. This version is the matrix-to-vector calculation on page 2; the original version carries the outer bond axes through the operation.

\newpage
\section*{5. MPS code: split and keep the updated MPS}
\setlength{\parskip}{3pt}
\setlength{\abovedisplayskip}{7pt}
\setlength{\belowdisplayskip}{7pt}
After CNOT, $\Theta_G$ contains the updated coefficients:
\[
(\Theta_G)_{0ij0}=(C_G)_{ij}.
\]
\subsection*{Arrange them as a matrix}
\begin{lstlisting}
left_bond = theta.shape[0]    # 1
right_bond = theta.shape[-1]  # 1
matrix = theta.reshape(left_bond * 2, 2 * right_bond)
\end{lstlisting}
\[
\underbrace{\Theta_G}_{(1,2,2,1)}
\xrightarrow{\text{reshape}}
\underbrace{C_G=\frac1{\sqrt2}\mtwo{1}{0}{0}{1}}_{\texttt{matrix},\ (2,2)}.
\]
In a longer chain, this reshape groups the left bond with the first spin into a row, and the second spin with the right bond into a column.

\subsection*{Factor into the new tensors}
\[
C_G=USV^\dagger,\qquad s_1=s_2=\frac1{\sqrt2}.
\]
\begin{lstlisting}
U, s, Vh = np.linalg.svd(matrix, full_matrices=False)
bond = max(1, int(np.count_nonzero(s > 1e-12)))
\end{lstlisting}
Both singular values survive the numerical threshold, so \texttt{bond = 2}.
\[
A_G^{[1]}=U,\qquad A_G^{[2]}=SV^\dagger.
\]
\begin{lstlisting}
A1 = U[:, :bond].reshape(left_bond, 2, bond)
A2 = (np.diag(s[:bond]) @ Vh[:bond, :]).reshape(
    bond, 2, right_bond
)
\end{lstlisting}
\textbf{The final MPS is now complete:} \texttt{A1} and \texttt{A2} have shapes $(1,2,2)$ and $(2,2,1)$. Keep these tensors for further gates or expectation values.

\subsection*{Optional: reconstruct only to verify}
\begin{lstlisting}
mps_after_CNOT = np.tensordot(
    A1, A2, axes=([-1], [0])
).reshape(-1)
print(np.linalg.norm(mps_after_CNOT - state_after_CNOT))
\end{lstlisting}
\[
A_G^{[1]}A_G^{[2]}=C_G
\xrightarrow{\text{flatten}}
\frac1{\sqrt2}\vfour{1}{0}{0}{1}.
\]
\end{document}